\chapter{역사적 배경}\label{ch:history}

이 장은 봄베가 어떤 문제를 풀기 위해, 누구에 의해, 어떤 순서로 만들어졌는지를 요약한다.
기계의 동작 원리는 \ref{ch:enigma}장과 \ref{ch:bombe}장에서 다룬다.

\section{에니그마}

에니그마는 글자를 누를 때마다 치환표가 바뀌는 로터식 암호기다. 1920년대에 상업용으로 나온 뒤 독일군이
채택했고, 군용 기종에는 글자 쌍을 맞바꾸는 플러그보드가 추가되었다~\cite{welchman1982hutsix,hinsley1993codebreakers}.
이 프로젝트가 다루는 기종은 세 가지다(표~\ref{tab:models}). 육군과 공군의 Enigma~I은 로터 다섯 개 중
세 개를 골라 쓰고, 해군의 M3는 로터를 여덟 개로 늘렸으며, U보트 부대의 M4는 네 번째 로터(그리스 로터)와
얇은 반사판을 더했다. 4로터 기종에서는 로터 순서 하나당 탐색할 위치가 $26^3$에서 $26^4$으로 26배
늘어난다~\cite{wenger1944history}.

\begin{table}[ht]
\centering
\caption{재현 대상 에니그마 기종. 로터 순서 수는 로터를 끼우는 순서의 가짓수이며 반사판 선택은 포함하지 않는다.}
\label{tab:models}
\begin{tabular}{llll}
\toprule
기종 & 로터 & 반사판 & 로터 순서 수 \\
\midrule
Enigma I (육군·공군) & I--V 중 3개 & B, C & $5\cdot4\cdot3 = 60$ \\
M3 (해군) & I--VIII 중 3개 & B, C & $8\cdot7\cdot6 = 336$ \\
M4 (해군 U보트) & I--VIII 중 3개 + Beta/Gamma & 얇은 B, 얇은 C & $336\cdot 2\cdot 2 = 1{,}344$ \\
\bottomrule
\end{tabular}
\end{table}

\section{폴란드의 선행 연구와 블레츨리 파크}

에니그마를 처음 수학적으로 공략한 것은 폴란드 암호국이다. 마리안 레예프스키는 1930년대 초 군용 에니그마의
로터 배선을 복원했고, 폴란드 팀은 당시 독일군의 지시자(indicator) 절차가 가진 약점을 이용하는 전기기계식
장치 ``봄바(bomba)''를 만들었다~\cite{rejewski1981polish}. 전쟁 직전인 1939년 여름 폴란드는 이 성과를
영국과 프랑스에 넘겼다~\cite{rejewski1981polish,welchman1982hutsix}.

영국의 암호 해독 기관인 정부암호학교(GC\&CS)는 블레츨리 파크에 자리를 잡았고, 앨런 튜링은 전쟁이 시작된
1939년 가을부터 그곳에서 에니그마를 맡았다~\cite{hodges1983turing}. 폴란드의 봄바는 지시자 절차에 의존했으므로
독일군이 절차를 바꾸면 쓸 수 없게 된다. 튜링은 절차가 아니라 \emph{추정 평문}(크립)에 기대는 새 기계를
설계했다. 전문 내용의 일부를 짐작할 수 있기만 하면 통하는 방법이다~\cite{hodges1983turing,copeland2004essential}.
이것이 봄베다.

\section{튜링의 봄베와 웰치먼의 대각선 보드}

튜링의 설계는 British Tabulating Machine(BTM)사가 제작했다. 첫 기계는 1940년 봄에 가동되었으나, 튜링의
원래 배선은 루프(폐회로)가 여러 개 있는 메뉴에서만 쓸모가 있었고 그런 크립은
드물었다~\cite{carter2010bombe,welchman1982hutsix}. 고든 웰치먼은 플러그보드가 대칭이라는 사실을
배선으로 강제하는 \term{대각선 보드}{diagonal board}를 고안했다. 대각선 보드가 없으면 동시 탐색에 루프가
세 개 이상 필요했지만, 있으면 루프가 하나거나 특수한 경우에는 없어도 된다~\cite{carter2010bombe}.
이후의 봄베는 모두 이 개량을 포함하며, 이 보고서에서 ``봄베''는 튜링--웰치먼 봄베를 가리킨다.

표준 3로터 봄베는 에니그마 등가 유닛(드럼 세 개 한 벌) 36개를 12개씩 세 뱅크로 배열했고, 드럼 108개와
지시 드럼 3개를 가졌다~\cite[5]{us6812report1944}. 폭 약 7피트, 높이 6피트 6인치, 무게 약 1톤의
기계다~\cite{ellsbury1998bombe,us6812report1944}. 정지가 없을 때 로터 순서 하나를 훑는 데 약 10.4분이
걸렸다~\cite{jackson-bh09,jackson-bh11}. 영국에서 제작된 봄베는 3로터와 4로터 기종을 합쳐
211대로 집계된다~\cite{jackson-bh14}.

그림~\ref{fig:bombe-front}은 문헌의 기술을 바탕으로 그린 이 기계의 모습이다. 앞면에는 드럼이 아홉 줄로
놓인다. 세 줄이 한 뱅크이고, 세로로 놓인 드럼 세 개가 스크램블러 한 벌, 곧 플러그보드 없는 에니그마 한
대다. 가운데 뱅크의 오른쪽에 정지 위치를 읽는 지시 드럼 세 개가 있다~\cite[5]{us6812report1944}.
뒷면은 메뉴를 배선하는 곳이다. 잭이 아홉 열로 놓이는데, 뱅크마다 대각선 보드의 잭 열, 커먼 잭 열,
스크램블러의 입출력 잭 열이 하나씩이다~\cite[59, 68]{us6812report1944}\cite{jackson-bh10}. 조작자는 26선
케이블로 스크램블러의 양 끝을 메뉴의 두 글자에 해당하는 잭에 꽂았다~\cite[8--10]{us6812report1944}.
반사판은 기계 왼쪽 측면의 플러그보드로 뱅크마다 설정했다~\cite[5, 39]{us6812report1944}.
각 부분이 하는 일은 \ref{ch:bombe}장에서 설명한다.

\begin{figure}[p]
\centering
% 영국 3로터 튜링-웰치먼 봄베의 앞면(위)과 뒷면 잭 배치(아래). 직접 그린 도해이며 사진을 옮긴 것이 아니다.
% 근거: docs/reference/bombe-mechanics.md 2, 5, 11절 (us6812report1944, ellsbury1998bombe, jackson-bh10).
% 단위: 1 = 1 ft. 근거가 있는 것: 뱅크 3개 x 유닛 12개 x 드럼 3개(108개), 지시 드럼 3개가 가운데 뱅크
% 오른쪽, 외형 약 7 ft x 6 ft 6 in, 드럼 지름 약 5 in, 뒷면 잭 열 9개의 용도.
% 어림해 그린 것: 드럼 사이 간격과 캐비닛 안에서의 정확한 위치, 지시 드럼의 배열 방향, 바퀴, 잭의 개수와 간격.
\begin{tikzpicture}[x=1.27cm,y=1.27cm,
    drum/.style={circle,draw=black!70,fill=clight,minimum size=5.3mm,inner sep=0pt,line width=0.4pt},
    lab/.style={font=\footnotesize,align=left,anchor=west},
    lead/.style={draw=cgray,line width=0.4pt}]
  % ---------------- 앞면 ----------------
  \draw[line width=0.9pt,fill=white] (0,0) rectangle (7,6.5);
  \foreach \x in {0.5,6.5} { \draw[fill=cgray!40] (\x,-0.12) circle (0.1); }
  % 드럼: 뱅크 3개 x 3줄 x 12열
  \foreach \b/\by in {0/4.55,1/2.8,2/1.05} {
    \fill[cbluebg!50] (0.33,\by-0.25) rectangle (6.02,\by+1.25);
    \foreach \r in {0,1,2} {
      \foreach \c in {0,...,11} {
        \node[drum] at (0.6+\c*0.47,\by+1.0-\r*0.5) {};
        \fill[black!70] (0.6+\c*0.47,\by+1.0-\r*0.5) circle (0.03);
      }
    }
  }
  % 스크램블러 한 벌 강조 (위 뱅크, 셋째 열)
  \draw[corange,line width=1.1pt,rounded corners=2pt] (1.32,4.32) rectangle (1.76,5.78);
  % 지시 드럼 (가운데 뱅크 오른쪽)
  \foreach \r in {0,1,2} { \node[drum,fill=cyellowbg] (ind\r) at (6.5,3.8-\r*0.5) {}; }
  % 치수
  \draw[<->,cgray] (0,6.8) -- node[above,font=\footnotesize,text=cgray] {폭 약 7 ft (2.1 m)} (7,6.8);
  \draw[<->,cgray] (-0.35,0) -- node[above,rotate=90,font=\footnotesize,text=cgray] {높이 약 6 ft 6 in (2.0 m)} (-0.35,6.5);
  % 라벨
  \draw[lead] (6.02,5.3) -- (7.3,5.9);
  \node[lab] at (7.3,5.9) {\textbf{뱅크}(bank) × 3\\유닛 12개씩, 모두 36개};
  \draw[lead,corange] (1.54,5.78) -- (1.9,6.15) -- (6.5,6.15) -- (7.3,4.95);
  \node[lab,text=corange] at (7.3,4.75) {\textbf{스크램블러 한 벌}(유닛)\\세로로 놓인 드럼 3개\\= 플러그보드 없는 에니그마};
  \draw[lead] (ind0.east) -- (7.3,3.5);
  \node[lab] at (7.3,3.5) {\textbf{지시 드럼}(indicator drums) 3개\\정지한 위치를 읽는다};
  \node[lab] at (7.3,2.1) {한 벌 안의 드럼 (위에서 아래로)\\위: 연속 회전 --- 왼쪽 로터\\가운데 --- 가운데 로터\\아래: 가장 느림 --- 오른쪽 로터};
  \node[lab] at (7.3,0.75) {드럼 108개, 지름 약 5 in\\깊이 약 2 ft, 무게 약 1톤};
  \node[font=\small\bfseries,anchor=west] at (0,7.35) {앞면};
  % ---------------- 뒷면 (잭 배치 도식) ----------------
  \begin{scope}[yshift=-6.6cm]
  \node[font=\small\bfseries,anchor=west] at (0,3.75) {뒷면: 메뉴를 배선하는 잭 (배치 도식)};
  \draw[line width=0.9pt,fill=white] (0,0) rectangle (7,3.4);
  \foreach \g in {0,1,2} {
    \draw[cgray,dashed] (0.25+\g*2.25,0.15) rectangle (2.25+\g*2.25,3.25);
    % 대각선 보드 열
    \foreach \j in {0,...,12} { \fill[cblue] (0.55+\g*2.25,0.4+\j*0.215) rectangle ++(0.22,0.12); }
    % commons 열
    \foreach \j in {0,...,7} { \fill[cgreen] (1.14+\g*2.25,0.5+\j*0.34) rectangle ++(0.22,0.2); }
    % 스크램블러 IN/OUT 열 (유닛당 한 쌍, 12쌍)
    \foreach \j in {0,...,11} {
      \fill[corange] (1.68+\g*2.25,0.4+\j*0.235) rectangle ++(0.13,0.13);
      \fill[corange!55] (1.84+\g*2.25,0.4+\j*0.235) rectangle ++(0.13,0.13);
    }
  }
  % 메뉴 케이블 예 (개념)
  \draw[black!75,line width=0.9pt] (1.75,2.8) .. controls (1.2,2.3) and (1.0,1.9) .. (0.66,1.55);
  \draw[black!75,line width=0.9pt] (1.91,2.8) .. controls (2.3,2.0) and (1.6,1.4) .. (1.25,1.28);
  \draw[black!75,line width=0.9pt] (1.25,1.6) .. controls (0.9,1.2) and (0.9,0.9) .. (0.66,0.68);
  \draw[lead,cblue] (0.66,3.0) -- (0.66,3.5) -- (7.3,3.1);
  \node[lab,text=cblue] at (7.3,3.1) {\textbf{대각선 보드}(diagonal board)\\글자마다 26접점 잭, 뱅크마다 한 벌};
  \draw[lead,cgreen] (5.75,2.0) -- (7.3,2.0);
  \node[lab,text=cgreen] at (7.3,2.0) {\textbf{커먼}(commons)\\한 글자에 선이 여럿 모일 때의 병렬 잭};
  \draw[lead,corange] (6.47,1.0) -- (7.3,0.9);
  \node[lab,text=corange] at (7.3,0.9) {\textbf{스크램블러 IN/OUT 잭}\\유닛마다 한 쌍. 26선 케이블로 잇는다};
  \end{scope}
\end{tikzpicture}

\caption{영국 3로터 봄베의 도해(직접 그린 것이며 사진을 옮긴 것이 아니다). 위는 앞면, 아래는 뒷면의 잭 배치다.
뱅크와 유닛과 드럼의 수, 지시 드럼의 위치, 드럼과 로터의 대응, 뒷면 잭 열의 용도는
1차 자료~\cite[5, 59, 68, 73]{us6812report1944}를, 외형 치수와 드럼 지름은~\cite{ellsbury1998bombe,us6812report1944}를
따랐다. 드럼 사이의 간격과 캐비닛 안에서의 정확한 위치, 지시 드럼의 배열 방향, 바퀴의 모양, 뒷면 잭의 개수와
간격, 세 묶음과 뱅크의 대응은 확인하지 못해 어림해 그렸다. 실제 드럼은 로터 종류에 따라 색이
달랐으나~\cite[9]{us6812report1944} 그림에서는 구분하지 않았다.}
\label{fig:bombe-front}
\end{figure}

\section{4로터 문제와 미국의 봄베}\label{sec:us-bombes}

1942년 독일 해군 U보트 부대가 M4를 도입하자 탐색 공간이 26배로 늘었다. 미 해군은 이를 ``26배 많은 봄베
또는 26배 빠른 봄베''가 필요한 문제로 보고 고속 봄베 개발을 승인했다(1942년 9월)~\cite{wenger1944history}.
National Cash Register(NCR)사의 조지프 데시가 설계를 맡았고~\cite{desch1942memo}, 튜링은 1942년 12월
데이턴의 NCR를 방문해 설계를 검토했다~\cite{turing1942ncr}. 미 해군 봄베는 4휠 에니그마 등가 유닛 16개를
가졌고, 정지 감지를 전자식으로 했으며, 정지한 위치에서 스테커를 자동으로 검사해 통과한 것만
인쇄했다~\cite{ekhall-hallenberg2018}. 양산기는 1943년 6월 초부터 가동되었다~\cite{wenger1944history}.

미 육군은 다른 길을 택했다. 벨 연구소가 만든 육군 봄베는 회전 드럼을 없애고 전화 교환기용 릴레이로
로터의 치환과 위치 진행을 구현했다~\cite{dturing2021army}. 영국에서도 3로터 봄베에 초고속 네 번째 휠을
덧붙인 4로터 변형들이 만들어졌다~\cite{jackson-bh12}. 세 계열의 차이는 표~\ref{tab:families}와 같다.

\begin{table}[ht]
\centering
\caption{봄베 세 계열의 비교. 출처: 영국~\cite{us6812report1944,jackson-bh10,jackson-bh11},
미 해군~\cite{ekhall-hallenberg2018,turing1942ncr,desch1942memo}, 미 육군~\cite{dturing2021army}.
미 해군의 구동 시간은 직접 열람하지 못한 문헌을 거친 간접 인용이다.}
\label{tab:families}
\small
\begin{tabularx}{\linewidth}{l>{\raggedright\arraybackslash}X>{\raggedright\arraybackslash}X>{\raggedright\arraybackslash}X}
\toprule
 & 영국 3로터 (BTM) & 미 해군 4로터 (NCR) & 미 육군 (벨 연구소) \\
\midrule
스크램블링 & 회전 드럼 & 회전 휠 & 고정 배선 + 릴레이 \\
위치 진행 & 모터, 웜 구동, 자리올림 기구 & 모터, 단별 기어열 & 펄스로 구동되는 릴레이 \\
유닛 수 & 36 (12 × 3뱅크) & 16 & 144 (전체) \\
메뉴 설정 & 26선 케이블 배선 & 회전 스위치 & 제어반 \\
정지 감지 & 감지 릴레이 & 전자식 & 정지 없이 기록 \\
정지 이후 & 수동 체킹 머신 & 자동 검사 뒤 인쇄 & 부착 장치로 자동 체킹 \\
1회 구동 & 10.4분 + 정지 처리 & 4휠 약 20분 & 7분 \\
\bottomrule
\end{tabularx}
\end{table}

\section{이 프로젝트가 재현하는 것}

이 프로젝트의 최종 목표는 \emph{기계식 실물}이다. 드럼이 실제로 돌고 전류가 배선을 흘러 정지하는 봄베를
만들며, 전시관에 두고 관람객이 직접 메뉴를 배선하고 구동해 볼 수 있는 수준을 설계 요구사항으로 삼는다
(전시 여부는 아직 정해지지 않았다).

전시의 구성은 같은 봄베를 세 가지 방식으로 만든 세 대다. 기어로 드럼을 돌리는 \emph{기계식}, 회전 부품을
릴레이로 바꾼 \emph{릴레이식}, 움직이는 부품이 없는 \emph{전자식}(FPGA)이다. 앞의 둘은 \ref{sec:us-bombes}절에서
본 대로 1940년대에 실제로 있었던 두 방식이고, 셋째는 오늘날의 방식이다. 세 대를 나란히 두고 같은 메뉴를
걸어 같은 위치에서 정지하는 것을 보이면, 같은 계산이 기계에서 전기기계로, 다시 전자로 옮겨 가며 얼마나
빨라졌는지가 드러난다. 따라서 FPGA 구현은 실물을 돕는 전자 장치이기도 하지만 그 자체가 세 번째 전시물이다
(그림~\ref{fig:three-machines}).

\begin{figure}[ht]
\centering
\begin{tikzpicture}[x=1cm,y=1cm]
  \node[blkg,minimum width=40mm,minimum height=12mm] (sw) at (5.2,2.6) {소프트웨어 기준 모델\\{\footnotesize 세 기계가 내야 할 정지 목록}};
  \node[blko,minimum width=38mm,minimum height=20mm] (g) at (0,0) {\textbf{기계식}\\{\footnotesize 기어로 도는 드럼과 브러시}\\{\footnotesize 영국 봄베의 방식}};
  \node[blko,minimum width=38mm,minimum height=20mm] (r) at (5.2,0) {\textbf{릴레이식}\\{\footnotesize 릴레이 접점망과 계수기}\\{\footnotesize 미 육군 봄베의 방식}};
  \node[blkb,minimum width=38mm,minimum height=20mm] (f) at (10.4,0) {\textbf{전자식 (FPGA)}\\{\footnotesize 움직이는 부품 없음}\\{\footnotesize 내부 상태를 표시 장치로}};
  \draw[->,line width=1pt] (g) -- (r); \draw[->,line width=1pt] (r) -- (f);
  \draw[->,dashed] (sw) -- (r); \draw[->,dashed] (sw.west) -| (g); \draw[->,dashed] (sw.east) -| (f);
  \node[note] at (5.2,-1.6) {같은 메뉴 카드 $\Rightarrow$ 같은 정지. 위치 하나를 시험하는 시간은 오른쪽으로 갈수록 짧아진다.};
  \node[note] at (0,-2.25) {기계};
  \node[note] at (5.2,-2.25) {전기기계};
  \node[note] at (10.4,-2.25) {전자};
\end{tikzpicture}
\caption{이 프로젝트가 계획하는 세 대의 전시물(계획을 나타낸 개념도이며 완성된 기계의 모습이 아니다).
점선은 기준 모델과의 결과 대조다.}
\label{fig:three-machines}
\end{figure}

소프트웨어 모델은 세 기계 모두의 기준 모델이다. 세 기계가 내야 할 답을 정의하고, 각 기계의 결과를 대조하는
근거가 된다. 이 보고서의 장 순서는 이 구성을 따른다. \ref{ch:software}장은 기준 모델을, \ref{ch:rtl}장은
전자식 구현을, \ref{ch:hardware}장은 기계식과 릴레이식 실물을 다룬다. 전자식을 먼저 다루는 것은 역사의
순서와는 반대지만 제작의 순서다. 기준 모델과 비트 단위로 대조할 수 있는 회로를 먼저 완성해 두면 그것이
실물의 구동 제어와 결과 대조에도 쓰인다.

``재현''에는 두 기준이 있다. 첫째는 기능이다. 암호문과 크립이 주어지면 로터 순서, 로터 위치,
플러그보드 배선 후보를 찾아내고 검증해 실제로 해독한다. 둘째는 구조다. 결과만 맞추는 탐색기가 아니라
실제 기계의 구성요소를 일대일로 옮긴다. 영국 3로터 봄베가 기본 대상이고, 4로터는 자료가 가장 충분한
미 해군형의 동작(자동 검사 포함)을 따른다. 실물 단계에서는 영국식 회전 드럼(기어식)과 미 육군식
릴레이 구성(릴레이식)을 둘 다 만든다. 미 육군 봄베가 로터 한 단을 ``고정 배선 + 순환 시프트''로
구현했다는 기록~\cite{dturing2021army}이 릴레이식 설계의 출발점이다.

이 보고서의 역사적 서술은 1차 자료(미 육군 6812 파견대의 봄베 운용 보고서~\cite{us6812report1944},
데시의 설계 메모~\cite{desch1942memo}, 튜링의 NCR 방문 보고~\cite{turing1942ncr})와 재현 프로젝트의
해설~\cite{jackson-bh10,ellsbury1998bombe,carter2010bombe}에 근거한다. 출처끼리 어긋나는 수치는
1차 자료를 따랐다.
